\onlyinsubfile{\setcounter{section}{0}}
\section{Introduction}
\label{sec:chapter2-intro}

\par There is much literature \parencite{Algan2010} arguing that trust matters for aggregate economic outcomes such as long-term growth. \cite{jbpglg2016} shows that individual economic performance is not simply increasing in trust, but is maximized at intermediate levels. This paper contributes to the literature by asking whether a similar nonlinearity appears for realized respondent-level investment returns in the Household Retirement Survey (HRS).

\par The core argument of this paper is that if there is a link between trust and economic performance, then a similar (and possibly stronger) link should exist between trust and performance measures with respect to investment behavior. Inherent in financial environments where borrowing and lending are mediated through contracting between borrowers and lenders regarding a delayed payoff, is the \textit{promise of future repayment}. If trust is too low, households may miss out on high-return opportunities. On the other hand, if trust is too high, investors may be too ambitious and get cheated more often. I argue that, if this \say{right amount of trust} mechanism results in a hump-shaped relationship between income and trust, then this fact about the nature of financial contracting should imply that the hump-shaped relationship is especially pronounced for returns. 

\par  Recent evidence using Norwegian administrative data documents persistent idiosyncratic return differences between households and scale dependence with wealth \cite{aflgdmlp20}. Related work using U.S. panel data also emphasizes a persistent return component and its relevance for inequality and life-cycle outcomes \cite{Daminato2024}. These findings motivate interest in documenting this persistent component and explaining it with observable household characteristics. In particular, a contribution of this paper would be to document this persistent component across respondents \textit{holding constant this hump-shaped relationship between trust and returns.}

\par Replicating the Norwegian measurement environment in the U.S. is difficult. U.S. researchers typically combine survey data with partial administrative information, rather than observe a full administrative panel of wealth flows at the population level \cite{Kennickell2017b}. In this paper, I use the publicly-available HRS RAND Longitundinal 2022 version of the data, which is unusually useful for this question because it contains a \textit{measure of trust} as well as the core ingredients needed to construct returns over time: \textit{capital income, asset values across waves, debt, and transaction-related information}.

\par I first verify that the moderate trust ranges are associated with the highest income levels in the HRS, both contemporaneously (2020) and for average income over the period (2002-2022). I then test for a comparable hump shape regarding trust and returns to net wealth under a number of specifications. Since trust is observed in one wave, I treat it as time-invariant in pooled and panel regression models.\footnote{\cite{Lesmeister2019} is an example of a paper which also has trust as a time-invariant regressor.} The main results of the paper is that I not only recover the hump-shaped relationship between trust and returns to net wealth, but also control for it while characterizing the persistent component of returns across respondents. 

\par The paper proceeds as follows. Section~\ref{sec:litrev} reviews the relevant literature on trust, household returns, and panel methods for time-invariant regressors. Section~\ref{sec:data} describes the construction of trust, income, wealth, portfolio shares, and return measures, and presents the main descriptive statistics using the RAND longitudinal HRS product. Section~\ref{sec:chapter2-results} presents the main empirical results, including the average-return and pooled specifications, and the panel evidence on persistent heterogeneity in returns. I also consider the distribution of trust responses compared to the estimated trust values which maximize returns net wealth and reports model-implied mean predicted returns by race relative to the sample average. Section~\ref{sec:exts} extends and repeats the analysis by controlling for risk exposure through portfolio shares and shows that the main trust results are robust to this adjustment. Section~\ref{sec:concl} concludes. The determinants of trust, the established hump-shaped relationship between trust and income in the HRS, and other results can be found in the Appendix~\ref{app:chapter2}.
